\section{Related Work}

\subsection{Agent Memory Systems}

Recent work has explored a range of designs for giving LLM agents persistent memory. MemGPT frames long-term memory as an operating-system-style resource-management problem: the LLM decides when to move information between a finite ``core memory'' (analogous to RAM) and a larger ``archival memory'' (analogous to disk), enabling persistent multi-turn interaction within a fixed context window~\cite{packer2023memgpt}. Letta extends this line of work into a broader framework for agent memory operating across sessions and tasks~\cite{packer2023memgpt}. Mem0 takes a more production-oriented approach: it extracts fine-grained memory facts from user interactions, stores them in a shared vector store, and retrieves the most relevant facts at query time, emphasizing scalability and ease of integration for personalized agents~\cite{mem02025}. Graphiti, developed in the Zep ecosystem, organizes memory as a temporal knowledge graph where entities and relationships evolve over time, aiming to preserve richer structural information than flat fact stores~\cite{rasmussen2025zep}. A-MEM introduces a more agentic perspective: the memory system itself decides which memories to create, maintain, or consolidate during ongoing interaction, rather than relying on fixed extraction rules~\cite{amem2025}.

These systems differ substantially in retrieval granularity (facts, nodes, or graph episodes), memory organization (flat, hierarchical, or graph-structured), and update behavior (append-only, overwrite, or merge). However, they share a common evaluation limitation: they are typically assessed through downstream task success, user-level memory recall, or general question answering. These metrics are valuable, but they do not isolate a narrower question that is central to our benchmark: when several semantically similar states of the same evolving object coexist in memory, can the system retrieve the one that matches the queried stage, rather than the most recent or most topically relevant one? TSSMRBench is designed to test exactly this capability. By constructing explicit versioned state trajectories and asking questions that require stage-specific retrieval, our benchmark provides a targeted diagnostic that complements the broader evaluations used in these systems.

\subsection{Long-Term Memory Benchmarks}

Several recent benchmarks have advanced the evaluation of persistent agent memory. LoCoMo constructs long multi-session conversational scenarios and tests whether a system can recall facts, summaries, and event details mentioned across sessions spanning hundreds of turns~\cite{locomo2024}. LongMemEval focuses on five memory-related capabilities: knowledge retention, reasoning, forethought, understanding of temporal ordering, and handling of conflicting information, and evaluates them over distributed long histories~\cite{longmemeval2024}. MemoryAgentBench examines a complementary setting: agents must reuse information from prior interactions to accomplish later tasks, testing whether memory actually improves downstream performance in a task-oriented pipeline~\cite{memoryagentbench2025}.

Taken together, these benchmarks have moved evaluation beyond single-turn QA and made persistent memory a serious empirical target. At the same time, they mainly treat memory as a collection of facts, dialogue details, or episodes to be recalled later. They do not usually build explicit versioned state trajectories for the same underlying object, and they rarely ask for an earlier but stage-correct state when a later state is also present. In other words, they test whether the system can \emph{remember}, but not whether it can retrieve the \emph{right version} of a memory that has changed over time. TSSMRBench complements this line of work by making this distinction explicit: each scenario contains a full evolution trajectory, and questions are designed so that the correct answer depends on a specific stage, not on the most recent or most similar memory.

\subsection{Temporal Reasoning and Knowledge Update Benchmarks}

A second neighboring family comes from temporal reasoning. Temporal databases and temporal query languages provided early formal tools for representing and querying time-dependent data~\cite{jensen1999temporal,snodgrass1995tsql2}. In NLP, TORQUE introduces a reading-comprehension dataset focused on event temporal ordering, testing whether models can determine which of two events happened first~\cite{ning2020torque}. TimeDial studies temporal commonsense in dialogue, probing whether models understand implicit temporal relations such as duration and typical event order~\cite{qin2021timedial}. TimeQA targets time-sensitive question answering over Wikipedia articles whose facts change over time, asking questions conditioned on a specific time period~\cite{timeqa2021}. ETR-QA provides a comprehensive evaluation of event temporal reasoning abilities, covering multiple reasoning types such as ordering, span estimation, and hypothetical counterfactuals~\cite{luo2025etrqa}. MusTQ introduces multi-step temporal reasoning over temporal knowledge graphs, requiring models to chain multiple temporal relationships~\cite{mustq2024}. Test of Time evaluates LLM temporal reasoning through both semantic understanding and time-arithmetic tasks, measuring whether models can perform explicit date computation~\cite{fatemi2024testoftime,time2025}. These benchmarks are important because they test whether models can reason about order, duration, and time-conditioned facts. However, they evaluate the \emph{reasoning} capabilities of the model itself rather than the retrieval capabilities of the memory system that feeds the model. In TSSMRBench, the question is not whether the model can order events given their descriptions, but whether the memory system can first retrieve the right version-state evidence from a pool of semantically similar candidates.

A third neighboring family comes from retrieval and evidence-based QA. HotpotQA emphasizes multi-hop evidence aggregation, requiring systems to combine information from multiple retrieved passages~\cite{yang2018hotpotqa}. NarrativeQA studies understanding over long narrative evidence, testing comprehension of plot-level events and relationships~\cite{kocisky2018narrativeqa}. Natural Questions evaluates open-domain retrieval and answer extraction from a large web corpus using real user queries~\cite{kwiatkowski2019naturalquestions}. MuSiQue focuses on compositional multi-hop questions that require integrating several pieces of retrieved evidence in a specific order~\cite{trivedi2022musique}. RAG and its later variants study how retrieved evidence can be incorporated into generation, and survey work now provides a broad overview of retrieval-augmented methods and their failure modes~\cite{gao2023rag,lewis2020rag}. These works sharpen our understanding of retrieval and evidence use, but they assume a static or unversioned document collection. They do not directly target stage-specific state disambiguation inside an evolving memory trajectory; that is, the retrieval target is a relevant passage, not a specific version of a changing fact.

Finally, knowledge-update settings raise a different but related issue: how a memory system should incorporate newer information and avoid outdated answers. This perspective is highly relevant, but its default assumption is often that newer information ought to dominate older information. TSSMRBench differs in an important way. Our goal is not simply to test whether a memory system prefers the latest state. Instead, we test whether the memory system can retrieve the \emph{correct} state for the queried stage, even when that state is not the most recent one. This distinction matters because in many real-world scenarios (legal rule tracking, policy revision history, project plan evolution) the relevant answer may refer to an intermediate state, not the current one.

\begin{table}[!t]
\centering
\scriptsize
\caption{Comparison of TSSMRBench with representative neighboring benchmarks. LongMem indicates long-term memory focus; Spec-state indicates explicit retrieval of the stage-correct state; Lang. reports the primary language; Scale reports coarse benchmark size.}
\label{tab:benchmark-comparison}
\setlength{\tabcolsep}{1.5pt}
\begin{tabularx}{\columnwidth}{@{}>{\raggedright\arraybackslash}p{0.28\columnwidth}>{\centering\arraybackslash}p{0.09\columnwidth}>{\centering\arraybackslash}p{0.11\columnwidth}>{\centering\arraybackslash}p{0.14\columnwidth}>{\centering\arraybackslash}p{0.08\columnwidth}>{\raggedright\arraybackslash}X@{}}
\toprule
Benchmark & \shortstack{Long\\Mem} & Temporal & \shortstack{Spec-\\state} & Lang. & Scale \\
\midrule
LoCoMo~\cite{locomo2024} & Yes & Partial & No & EN & 1,986 QA \\
LongMemEval~\cite{longmemeval2024} & Yes & Partial & No & EN & 500 QA \\
\shortstack[l]{MemoryAgentBench\\\cite{memoryagentbench2025}} & Yes & Partial & No & EN & Multi-source suite \\
TimeQA~\cite{timeqa2021} & No & Yes & No & EN & $\approx$40K QA \\
MusTQ~\cite{mustq2024} & No & Yes & No & EN & $\approx$666K QA \\
\midrule
TSSMRBench & Yes & Yes & Yes & EN & 300 repositories / 900 QA \\
\bottomrule
\end{tabularx}
\end{table}
